\section{Problem Description}
\label{sec:problem_description}

\subsection{Physical Domain and Flow Conditions}

The objective of this project is to obtain a numerical solution
to the two-dimensional energy equation for laminar, incompressible
flow through a rectangular channel. The computational domain has
a length of $L_x=5$ and a height of $L_y=1$. A uniform,
cell-centered grid is used to discretize the domain.

The flow is assumed to be fully developed, with a prescribed
velocity profile in the streamwise direction and no velocity
component in the transverse direction. The dimensionless velocity
field is defined as
\begin{equation}
u(y)=6\overline{u}\,y(1-y),
\qquad
v(y)=0,
\label{eq:velocity_profile}
\end{equation}
where $\overline{u}$ is the mean streamwise velocity. For
$\overline{u}=3$, the velocity profile becomes
\begin{equation}
\boxed{
u(y)=18y(1-y),
\qquad v(y)=0.
}
\label{eq:velocity_profile_expanded}
\end{equation}

The nominal flow parameters are
\begin{equation}
\mathrm{Re}=25,
\qquad
\mathrm{Pr}=0.7,
\qquad
\mathrm{Ec}=0.1,
\label{eq:nominal_parameters}
\end{equation}
where $\mathrm{Re}$, $\mathrm{Pr}$, and $\mathrm{Ec}$ denote
the Reynolds, Prandtl, and Eckert numbers, respectively.

\subsection{Boundary and Initial Conditions}

The temperature field is subject to prescribed thermal
conditions at the channel boundaries. The lower and upper
walls are maintained at constant dimensionless temperatures
of zero and one, respectively:
\begin{align}
T(x,0,t)&=0,
\label{eq:bottom_wall_condition}
\\
T(x,1,t)&=1.
\label{eq:top_wall_condition}
\end{align}

At the inlet, the temperature is prescribed using the
linear profile
\begin{equation}
T(0,y,t)=T_{\mathrm{in}}(y)=y.
\label{eq:inlet_temperature_condition}
\end{equation}

At the outlet, a zero streamwise temperature gradient is
applied:
\begin{equation}
\left.\frac{\partial T}{\partial x}\right|_{x=L_x}=0.
\label{eq:outlet_temperature_condition}
\end{equation}
This condition represents zero conductive heat flux in the
streamwise direction at the outlet.

The initial temperature distribution throughout the domain
is specified as
\begin{equation}
\boxed{
T(x,y,0)=y.
}
\label{eq:initial_temperature}
\end{equation}

Thus, the initial temperature profile matches the prescribed
inlet profile and varies linearly between the lower and upper
wall temperatures. Starting from this initial condition, the
temperature field evolves under the combined effects of
advection, thermal diffusion, and viscous dissipation until
a steady-state solution is obtained.